% Manuscript styling
\usepackage{upgreek}
\captionsetup{font=singlespacing,justification=justified}

% Table formatting
\usepackage{longtable}
\usepackage{lscape}
% \usepackage[counterclockwise]{rotating}   % Landscape page setup for large tables
\usepackage{multirow}		% Table styling
\usepackage{tabularx}		% Control Column width
\usepackage[flushleft]{threeparttable}	% Allows for three part tables with a specified notes section
\usepackage{threeparttablex}            % Lets threeparttable work with longtable

% Create new environments so endfloat can handle them
% \newenvironment{ltable}
%   {\begin{landscape}\centering\begin{threeparttable}}
%   {\end{threeparttable}\end{landscape}}
\newenvironment{lltable}{\begin{landscape}\centering\begin{ThreePartTable}}{\end{ThreePartTable}\end{landscape}}

% Enables adjusting longtable caption width to table width
% Solution found at http://golatex.de/longtable-mit-caption-so-breit-wie-die-tabelle-t15767.html
\makeatletter
\newcommand\LastLTentrywidth{1em}
\newlength\longtablewidth
\setlength{\longtablewidth}{1in}
\newcommand{\getlongtablewidth}{\begingroup \ifcsname LT@\roman{LT@tables}\endcsname \global\longtablewidth=0pt \renewcommand{\LT@entry}[2]{\global\advance\longtablewidth by ##2\relax\gdef\LastLTentrywidth{##2}}\@nameuse{LT@\roman{LT@tables}} \fi \endgroup}

% \setlength{\parindent}{0.5in}
% \setlength{\parskip}{0pt plus 0pt minus 0pt}

% Overwrite redefinition of paragraph and subparagraph by the default LaTeX template
% See https://github.com/crsh/papaja/issues/292
\makeatletter
\renewcommand{\paragraph}{\@startsection{paragraph}{4}{\parindent}%
  {0\baselineskip \@plus 0.2ex \@minus 0.2ex}%
  {-1em}%
  {\normalfont\normalsize\bfseries\itshape\typesectitle}}

\renewcommand{\subparagraph}[1]{\@startsection{subparagraph}{5}{1em}%
  {0\baselineskip \@plus 0.2ex \@minus 0.2ex}%
  {-\z@\relax}%
  {\normalfont\normalsize\itshape\hspace{\parindent}{#1}\textit{\addperi}}{\relax}}
\makeatother

\makeatletter
\usepackage{etoolbox}
\patchcmd{\maketitle}
  {\section{\normalfont\normalsize\abstractname}}
  {\section*{\normalfont\normalsize\abstractname}}
  {}{\typeout{Failed to patch abstract.}}
\patchcmd{\maketitle}
  {\section{\protect\normalfont{\@title}}}
  {\section*{\protect\normalfont{\@title}}}
  {}{\typeout{Failed to patch title.}}
\makeatother

\usepackage{xpatch}
\makeatletter
\xapptocmd\appendix
  {\xapptocmd\section
    {\addcontentsline{toc}{section}{\appendixname\ifoneappendix\else~\theappendix\fi: #1}}
    {}{\InnerPatchFailed}%
  }
{}{\PatchFailed}
\makeatother

% ============================================
% 中文本地化配置 (papaja-zh)
% ============================================

% 加载 ctex 包（中文支持）
\usepackage{ctex}

% 覆盖 apa6 的英文标签
% 使用 \AtBeginDocument 确保在 english 选项初始化之后执行
\AtBeginDocument{
  \renewcommand{\abstractname}{摘要}
  \renewcommand{\keywordname}{关键词}
  \renewcommand{\rheadname}{短标题}
  \renewcommand{\acksname}{注}
  \renewcommand{\figurename}{图}
  \renewcommand{\tablename}{表}
}

% 阻止 APA6 man 模式下摘要后重复显示标题（消除空白页）
\makeatletter
\def\def@donotrepeattitle{}
\makeatother

% 修改 author note 格式："注："后直接加文字，不要换行居中
\makeatletter
\patchcmd{\maketitle}
  {\begin{center}%
            \acksname%
         \end{center}}
  {\noindent\textbf{\acksname：}\ignorespaces}
  {}{\typeout{Failed to patch author note title.}}
\patchcmd{\maketitle}
  {\indent\par\@acks}
  {\@acks}
  {}{\typeout{Failed to patch author note content.}}
\makeatother

% ============================================
% 双语支持配置 (papaja-zh)
% ============================================

% 定义存储英文内容的宏
\makeatletter
\newcommand{\@titleen}{}
\newcommand{\@authoren}{}
\newcommand{\@affiliationen}{}
\newcommand{\@abstracten}{}
\newcommand{\@keywordsen}{}

% 提供用户可调用的设置命令
\newcommand{\titleen}[1]{\gdef\@titleen{#1}}
\newcommand{\authoren}[1]{\gdef\@authoren{#1}}
\newcommand{\affiliationen}[1]{\gdef\@affiliationen{#1}}
\newcommand{\abstracten}[1]{\gdef\@abstracten{#1}}
\newcommand{\keywordsen}[1]{\gdef\@keywordsen{#1}}
\makeatother

% 重定义 \maketitle，在中文摘要后追加英文摘要页
% 注意：英文内容的宏定义由 Lua 过滤器生成，经 post_processor 移动到 preamble
\makeatletter
\let\oldmaketitle\maketitle
\renewcommand{\maketitle}{%
  \oldmaketitle%
  % 仅当存在英文摘要时才插入英文页
  \ifx\@abstracten\@empty\else%
    \newpage%
    % 英文标题（如有）
    \ifx\@titleen\@empty\else%
      \begin{center}%
        \textbf{\@titleen}%
      \end{center}%
      \vspace{0.25cm}%
    \fi%
    % 英文作者（如有）
    \ifx\@authoren\@empty\else%
      \begin{center}%
        \@authoren%
      \end{center}%
      \vspace{0.1cm}%
    \fi%
    % 英文机构（如有）
    \ifx\@affiliationen\@empty\else%
      \begin{center}%
        \@affiliationen%
      \end{center}%
      \vspace{0.25cm}%
    \fi%
    % 英文摘要
    {\section*{\normalfont\normalsize Abstract}}%
    \begin{quote}%
      \@abstracten%
    \end{quote}%
    % 英文关键词
    \ifx\@keywordsen\@empty\else%
      \indent\textit{Keywords:} \@keywordsen%
    \fi%
    \newpage%
  \fi%
}
\makeatother
